\begin{frame}{자주 하는 실수 (3): 임계값 0.5를 ``모델 성능''으로 보기}
  \begin{itemize}
    \item ``정확도가 왜 75\%야?''라고 물었을 때, 답은
      ``모델이 나쁘다''가 아니라
      \kb{``임계값 0.5가 이 문제의 비용 구조에 맞지
      않는다''}일 수 있다.
    \item 2.1절 선형회귀의 ``예측값''이 숫자 \textbf{자체}가
      최종 답인데, 로지스틱회귀의 ``예측 확률''은
      \textbf{임계값을 적용하기 전의 중간 값}.
  \end{itemize}
  \vskip0.8em
  \begin{block}{정확한 순서}
    ``모델이 나쁘다''고 하기 전에:
    \begin{enumerate}
      \item \kb{PR-AUC / ROC-AUC} 로 순위 능력을 \textbf{먼저} 확인,
      \item \kb{비용 구조에 맞는 임계값}을 sweeping해서
        ``이 임계값에서 나쁘다''를 확인,
      \item 그때 ``모델이 나쁘다''고 결론.
    \end{enumerate}
\vskip0.4em
순서가 반 거꾸면 ``모델을 갈아엎어야 하나''에서
    ``임계값을 0.14로 바꾸면 된다''로 바뀌지 \emph{않}는다.
  \end{block}
\end{frame}
